\documentclass[12pt,a4paper]{article}

% Fonts — LuaLaTeX
\usepackage{fontspec}
\usepackage{unicode-math}
\setmainfont{Linux Libertine O}
\setmathfont{Latin Modern Math}
\newfontfamily\igprimfont{Everson Mono}
\newfontfamily\shavfont[Scale=1.0]{Trabajo}
\setmonofont[Scale=0.85]{DejaVu Sans Mono}
\newcommand{\igfont}{\shavfont}

% Shavian Unicode block (U+10450–U+1047F) → Trabajo automatically
% Works in text mode, inside \text{} in math, and in table cells.
\usepackage{newunicodechar}
\newunicodechar{𐑐}{\ifmmode\text{{\shavfont 𐑐}}\else{{\shavfont 𐑐}}\fi}
\newunicodechar{𐑑}{\ifmmode\text{{\shavfont 𐑑}}\else{{\shavfont 𐑑}}\fi}
\newunicodechar{𐑒}{\ifmmode\text{{\shavfont 𐑒}}\else{{\shavfont 𐑒}}\fi}
\newunicodechar{𐑓}{\ifmmode\text{{\shavfont 𐑓}}\else{{\shavfont 𐑓}}\fi}
\newunicodechar{𐑔}{\ifmmode\text{{\shavfont 𐑔}}\else{{\shavfont 𐑔}}\fi}
\newunicodechar{𐑕}{\ifmmode\text{{\shavfont 𐑕}}\else{{\shavfont 𐑕}}\fi}
\newunicodechar{𐑖}{\ifmmode\text{{\shavfont 𐑖}}\else{{\shavfont 𐑖}}\fi}
\newunicodechar{𐑗}{\ifmmode\text{{\shavfont 𐑗}}\else{{\shavfont 𐑗}}\fi}
\newunicodechar{𐑘}{\ifmmode\text{{\shavfont 𐑘}}\else{{\shavfont 𐑘}}\fi}
\newunicodechar{𐑙}{\ifmmode\text{{\shavfont 𐑙}}\else{{\shavfont 𐑙}}\fi}
\newunicodechar{𐑚}{\ifmmode\text{{\shavfont 𐑚}}\else{{\shavfont 𐑚}}\fi}
\newunicodechar{𐑛}{\ifmmode\text{{\shavfont 𐑛}}\else{{\shavfont 𐑛}}\fi}
\newunicodechar{𐑜}{\ifmmode\text{{\shavfont 𐑜}}\else{{\shavfont 𐑜}}\fi}
\newunicodechar{𐑝}{\ifmmode\text{{\shavfont 𐑝}}\else{{\shavfont 𐑝}}\fi}
\newunicodechar{𐑞}{\ifmmode\text{{\shavfont 𐑞}}\else{{\shavfont 𐑞}}\fi}
\newunicodechar{𐑟}{\ifmmode\text{{\shavfont 𐑟}}\else{{\shavfont 𐑟}}\fi}
\newunicodechar{𐑠}{\ifmmode\text{{\shavfont 𐑠}}\else{{\shavfont 𐑠}}\fi}
\newunicodechar{𐑡}{\ifmmode\text{{\shavfont 𐑡}}\else{{\shavfont 𐑡}}\fi}
\newunicodechar{𐑢}{\ifmmode\text{{\shavfont 𐑢}}\else{{\shavfont 𐑢}}\fi}
\newunicodechar{𐑣}{\ifmmode\text{{\shavfont 𐑣}}\else{{\shavfont 𐑣}}\fi}
\newunicodechar{𐑤}{\ifmmode\text{{\shavfont 𐑤}}\else{{\shavfont 𐑤}}\fi}
\newunicodechar{𐑥}{\ifmmode\text{{\shavfont 𐑥}}\else{{\shavfont 𐑥}}\fi}
\newunicodechar{𐑦}{\ifmmode\text{{\shavfont 𐑦}}\else{{\shavfont 𐑦}}\fi}
\newunicodechar{𐑧}{\ifmmode\text{{\shavfont 𐑧}}\else{{\shavfont 𐑧}}\fi}
\newunicodechar{𐑨}{\ifmmode\text{{\shavfont 𐑨}}\else{{\shavfont 𐑨}}\fi}
\newunicodechar{𐑩}{\ifmmode\text{{\shavfont 𐑩}}\else{{\shavfont 𐑩}}\fi}
\newunicodechar{𐑪}{\ifmmode\text{{\shavfont 𐑪}}\else{{\shavfont 𐑪}}\fi}
\newunicodechar{𐑫}{\ifmmode\text{{\shavfont 𐑫}}\else{{\shavfont 𐑫}}\fi}
\newunicodechar{𐑬}{\ifmmode\text{{\shavfont 𐑬}}\else{{\shavfont 𐑬}}\fi}
\newunicodechar{𐑭}{\ifmmode\text{{\shavfont 𐑭}}\else{{\shavfont 𐑭}}\fi}
\newunicodechar{𐑮}{\ifmmode\text{{\shavfont 𐑮}}\else{{\shavfont 𐑮}}\fi}
\newunicodechar{𐑯}{\ifmmode\text{{\shavfont 𐑯}}\else{{\shavfont 𐑯}}\fi}
\newunicodechar{𐑰}{\ifmmode\text{{\shavfont 𐑰}}\else{{\shavfont 𐑰}}\fi}
\newunicodechar{𐑱}{\ifmmode\text{{\shavfont 𐑱}}\else{{\shavfont 𐑱}}\fi}
\newunicodechar{𐑲}{\ifmmode\text{{\shavfont 𐑲}}\else{{\shavfont 𐑲}}\fi}
\newunicodechar{𐑳}{\ifmmode\text{{\shavfont 𐑳}}\else{{\shavfont 𐑳}}\fi}
\newunicodechar{𐑴}{\ifmmode\text{{\shavfont 𐑴}}\else{{\shavfont 𐑴}}\fi}
\newunicodechar{𐑵}{\ifmmode\text{{\shavfont 𐑵}}\else{{\shavfont 𐑵}}\fi}
\newunicodechar{𐑶}{\ifmmode\text{{\shavfont 𐑶}}\else{{\shavfont 𐑶}}\fi}
\newunicodechar{𐑷}{\ifmmode\text{{\shavfont 𐑷}}\else{{\shavfont 𐑷}}\fi}
\newunicodechar{𐑸}{\ifmmode\text{{\shavfont 𐑸}}\else{{\shavfont 𐑸}}\fi}
\newunicodechar{𐑹}{\ifmmode\text{{\shavfont 𐑹}}\else{{\shavfont 𐑹}}\fi}
\newunicodechar{𐑺}{\ifmmode\text{{\shavfont 𐑺}}\else{{\shavfont 𐑺}}\fi}
\newunicodechar{𐑻}{\ifmmode\text{{\shavfont 𐑻}}\else{{\shavfont 𐑻}}\fi}
\newunicodechar{𐑼}{\ifmmode\text{{\shavfont 𐑼}}\else{{\shavfont 𐑼}}\fi}
\newunicodechar{𐑽}{\ifmmode\text{{\shavfont 𐑽}}\else{{\shavfont 𐑽}}\fi}
\newunicodechar{𐑾}{\ifmmode\text{{\shavfont 𐑾}}\else{{\shavfont 𐑾}}\fi}
\newunicodechar{𐑿}{\ifmmode\text{{\shavfont 𐑿}}\else{{\shavfont 𐑿}}\fi}

% Page layout
\usepackage[top=1in, bottom=1in, left=1in, right=1in]{geometry}
\usepackage{microtype}
\usepackage{parskip}

% Language
\usepackage[english]{babel}

% Headers
\usepackage{fancyhdr}
\pagestyle{fancy}
\fancyhf{}
\fancyhead[L]{\leftmark}
\fancyhead[R]{\thepage}
\fancyfoot[C]{\thepage}
\renewcommand{\headrulewidth}{0.4pt}
\renewcommand{\footrulewidth}{0pt}

% Section styling — fully unnumbered; pipeline never auto-numbers
\usepackage{titlesec}
\setcounter{secnumdepth}{0}
\titleformat{\section}{\Large\bfseries}{}{0em}{}
\titleformat{\subsection}{\large\bfseries}{}{0em}{}
\titleformat{\subsubsection}{\normalsize\bfseries}{}{0em}{}

% Essential packages
\usepackage{graphicx}
\graphicspath{{/home/mrnob0dy666/imsgct/Smaragdine_Synthetica/manuscripts/alchemical_theory_of_math/}{/home/mrnob0dy666/imsgct/Smaragdine_Synthetica/manuscripts/alchemical_theory_of_math/images/}}
\setkeys{Gin}{width=0.48\linewidth,height=0.32\textheight,keepaspectratio}
\usepackage[table]{xcolor}
\usepackage{booktabs}
\usepackage{array}
\usepackage{tabularx}
\usepackage{longtable}
\usepackage{amsmath}
% unicode-math subsumes amssymb — do not load amssymb alongside it
\usepackage{float}
\usepackage[normalem]{ulem}
\renewcommand{\arraystretch}{1.3}

% Cross-references
\usepackage{hyperref}
\usepackage{cleveref}
\hypersetup{colorlinks=true, linkcolor=blue, urlcolor=cyan, citecolor=teal}

% Custom commands
\newcommand{\shav}[1]{{\shavfont #1}}
\newcommand{\heb}[1]{{\igprimfont #1}}
\newcommand{\tupleaddr}[1]{$\langle #1 \rangle$}

% Shorthand primitive commands — redefined for IG document use
\renewcommand{\H}{Ħ}
\newcommand{\K}{Ç}
\newcommand{\G}{Γ}
\newcommand{\g}{ɢ}
\newcommand{\Th}{Þ}
\newcommand{\D}{Ð}
\newcommand{\R}{Ř}
\newcommand{\f}{ƒ}

% Imscription tuple box
\usepackage{tcolorbox}
\tcbuselibrary{skins,breakable}
\newtcolorbox{imscriptionbox}{enhanced, colback=blue!5, colframe=blue!75!black, fontupper=\large, center upper, boxrule=1pt, arc=4pt}

% PDF metadata
\hypersetup{
  pdftitle={01 Conventional Decomposition},
  pdfkeywords={Imscribing Grammar},
  pdfauthor={Lando Mills},
}

\title{01 Conventional Decomposition}
\date{\today}
\author{Lando Mills}

\begin{document}
\maketitle
\hypertarget{conventional-decomposition-of-the-12-primitive-crystal}{%
\section{Conventional Decomposition of the 12-Primitive
Crystal}\label{conventional-decomposition-of-the-12-primitive-crystal}}

Phase 1 of the Alchemical Theory of Mathematics. This document stands
alone as a reference: it pushes each per-primitive CL8NK fragment down
into two layers, a first-order / ZFC unfolding and the named
conventional structure it denotes. No alchemical claim is made here. The
thesis that the primitive order is a derivation order is argued in the
companion manuscript and rests on this table.

\hypertarget{sources-and-provenance}{%
\subsection{Sources and provenance}\label{sources-and-provenance}}

Every fragment below is quoted from \texttt{CL8NK\_FORMULAE} in
\texttt{imscribing\_grammar/navigators/cl8nk\_navigator.py} (the
per-primitive formula map, sourced in turn from
\texttt{IG\_catalog.json}). Axis names and value orderings are the
authoritative ones from \texttt{Smaragdine\_Synthetica/lean/Core.lean}
(v0.5.69), whose constructor order defines the ordinals. Nothing here is
hand derived; each row is a fragment as emitted, annotated with its
conventional reading.

\textbf{The proximity labels are L8-relative.} \texttt{apex},
\texttt{close} and \texttt{distant} are the navigator's distance from
the CLINK L8 reference, not an absolute ranking of the values. This
mattered little while L8 was the terminal layer of the chain. It matters
now: CLINK L9 was born 2026-07-13 and sits above L8 at O\_∞⁺, with its
own reference typing and its own navigator
(\texttt{cl9nk\_navigator.py}). Several values this table marks
\texttt{distant} are apex in L9's register. Read the labels as
coordinates in a dialect, not as a verdict on a value.

\hypertarget{dialect-reconciliation-navigator-glyphs-to-core-axes}{%
\subsection{Dialect reconciliation (navigator glyphs to Core
axes)}\label{dialect-reconciliation-navigator-glyphs-to-core-axes}}

The navigator labels the twelve axes with one glyph alphabet; Core.lean
labels them with another. They are the same twelve axes in the same
tuple positions. The correspondence below is fixed by fragment content
and confirmed by family cardinality: each pair shares the same value
count (𝓕₃ = 3, 𝓕₄ = 4, 𝓕₅ = 5).

\begin{longtable}[]{@{}rclcl@{}}
\toprule
Pos & Navigator & Core axis & Family & Conventional axis
meaning\tabularnewline
\midrule
\endhead
1 & Ð & D Dimensionality & 𝓕₄ & recursive nesting depth of the state
space\tabularnewline
2 & Þ & T Topology & 𝓕₅ & connectivity pattern of the state
graph\tabularnewline
3 & Ř & R Relational Mode & 𝓕₄ & direction and duality of the binding
relation\tabularnewline
4 & Φ & P Parity / Symmetry & 𝓕₅ & symmetry group fixing the
object\tabularnewline
5 & ƒ & F Fidelity & 𝓕₃ & information loss of the read/write
map\tabularnewline
6 & Ç & K Kinetic Character & 𝓕₅ & potential held, and whether its gate
is open\tabularnewline
7 & Γ & G Scope / Granularity & 𝓕₃ & subset-size regime of
correlations\tabularnewline
8 & ɢ & Γ Interaction Grammar & 𝓕₄ & logical connective of the
coupling\tabularnewline
9 & ⊙ & Φ Criticality & 𝓕₅ & analytic character of the fixed
point\tabularnewline
10 & Ħ & H Chirality & 𝓕₄ & handedness: the direction a held potential
can go\tabularnewline
11 & Σ & S Stoichiometry & 𝓕₃ & matching cardinality of the two
sides\tabularnewline
12 & Ω & Ω Topological Protection & 𝓕₄ & homotopy invariant guarding the
state\tabularnewline
\bottomrule
\end{longtable}

Notation. \texttt{μ∘δ\ =\ id} is the special Frobenius condition: a
comonoid split \texttt{δ} followed by a monoid fusion \texttt{μ} returns
the identity, i.e.~\texttt{δ} is a section and \texttt{μ} a retraction
of it (a split idempotent that is in fact the identity on the object).
\texttt{V\ =\ L(x)} is the constructible-universe predicate.
\texttt{rank} is the von Neumann set-theoretic rank. \texttt{S} denotes
a shift/mirror successor operator.

\begin{center}\rule{0.5\linewidth}{0.5pt}\end{center}

\hypertarget{axis-1.-dimensionality-d-navigator-uxf0-ux1d4d5ux2084}{%
\subsection{Axis 1. Dimensionality (D, navigator Ð,
𝓕₄)}\label{axis-1.-dimensionality-d-navigator-uxf0-ux1d4d5ux2084}}

Conventional meaning: recursive nesting depth of the state space, from a
flat sheet up to a boundary-bulk imscriptive fixed point.

\begin{itemize}
\tightlist
\item
  \textbf{𐑦 {[}HOLOGRAPHIC\_STATE{]}, apex.} Fragment
  \texttt{V\ =\ L(x)\ ∧\ selfmodel(x)\ ∧\ x\ ∈\ V}.

  \begin{itemize}
  \tightlist
  \item
    FO: \texttt{V\ =\ L(x)\ ∧\ (∃m\ ∈\ x)(m\ codes\ x)\ ∧\ x\ ∈\ V}. The
    object equals its own constructible closure, contains a code of
    itself, and lies in the universe it generates.
  \item
    Named: a \textbf{reflexive / self-modeling object}; boundary-bulk
    imscriptive correspondence (the boundary datum determines the bulk).
    Fixed point of the constructibility operator, cf.~Gödel \texttt{L}
    and a self-referential coding (diagonal lemma).
  \end{itemize}
\item
  \textbf{𐑼, close.} Fragment
  \texttt{∀n∃y(y\ ∈\ x\ ∧\ rank(y)\ \textgreater{}\ n)}.

  \begin{itemize}
  \tightlist
  \item
    FO: as written; \texttt{x} has members of unbounded rank.
  \item
    Named: \textbf{rank-unbounded set}; not set-like of bounded height,
    i.e.~a proper class-like or limit-rank object. Infinite-dimensional
    in the rank filtration.
  \end{itemize}
\item
  \textbf{𐑨, distant.} Fragment \texttt{dim(x)\ =\ 2\ ∧\ sur(x)}.

  \begin{itemize}
  \tightlist
  \item
    FO: \texttt{dim(x)\ =\ 2\ ∧\ surface(x)}.
  \item
    Named: a \textbf{2-manifold / surface}; the flat sheet, homological
    dimension 2.
  \end{itemize}
\item
  \textbf{𐑛, distant.} Fragment \texttt{dim(x)\ =\ 0\ ∧\ fin(x)}.

  \begin{itemize}
  \tightlist
  \item
    FO: \texttt{dim(x)\ =\ 0\ ∧\ finite(x)}.
  \item
    Named: a \textbf{0-dimensional finite set}; discrete points,
    hereditarily finite.
  \end{itemize}
\end{itemize}

\begin{center}\rule{0.5\linewidth}{0.5pt}\end{center}

\hypertarget{axis-2.-topology-t-navigator-uxfe-ux1d4d5ux2085}{%
\subsection{Axis 2. Topology (T, navigator Þ,
𝓕₅)}\label{axis-2.-topology-t-navigator-uxfe-ux1d4d5ux2085}}

Conventional meaning: connectivity pattern of the state graph.

\begin{itemize}
\tightlist
\item
  \textbf{𐑸 {[}HOLOBOUND{]}, apex.} Fragment
  \texttt{bound\_⊙(a,\ f)\ ∧\ Refl(a,\ f)\ ∧\ holo(x,\ a)}.

  \begin{itemize}
  \tightlist
  \item
    FO: \texttt{boundary(a)\ ∧\ reflects(a,\ f)\ ∧\ imscribes(x,\ a)};
    the boundary \texttt{a} reflects the field \texttt{f} and losslessly
    imscribes the bulk \texttt{x}.
  \item
    Named: a \textbf{boundary-bulk correspondence} (bounded reflection);
    the bulk is reconstructible from boundary data.
  \end{itemize}
\item
  \textbf{𐑥, close.} Fragment \texttt{cross(x,\ y)\ ∧\ ¬\ meet(x,\ y)}.

  \begin{itemize}
  \tightlist
  \item
    FO: \texttt{cross(x,\ y)\ ∧\ ¬∃z(z\ ≤\ x\ ∧\ z\ ≤\ y)}.
  \item
    Named: a \textbf{bowtie / transversal crossing} with empty meet; two
    cycles sharing a node but no common lower bound (a figure-eight,
    bifurcation point).
  \end{itemize}
\item
  \textbf{𐑶, distant.} Fragment \texttt{x\ ⊠\ y\ ∧\ irreducible(x,\ y)}.

  \begin{itemize}
  \tightlist
  \item
    FO: \texttt{x\ ⊠\ y\ ∧\ irreducible(x,\ y)}.
  \item
    Named: an \textbf{irreducible tensor / box product}; no nontrivial
    factorization.
  \end{itemize}
\item
  \textbf{𐑡, distant.} Fragment \texttt{graph(x)\ ∧\ branch(x)}.

  \begin{itemize}
  \tightlist
  \item
    FO: \texttt{graph(x)\ ∧\ (∃v)\ deg(v)\ ≥\ 3}.
  \item
    Named: a \textbf{branching graph / tree with vertices of degree ≥
    3}.
  \end{itemize}
\item
  \textbf{𐑰, distant.} Fragment \texttt{x\ ⊆\ y\ ∧\ cont(y)}.

  \begin{itemize}
  \tightlist
  \item
    FO: \texttt{x\ ⊆\ y\ ∧\ continuous(y)}.
  \item
    Named: \textbf{nested inclusion into a continuum}; hierarchical
    containment.
  \end{itemize}
\end{itemize}

\begin{center}\rule{0.5\linewidth}{0.5pt}\end{center}

\hypertarget{axis-3.-relational-mode-r-navigator-ux159-ux1d4d5ux2084}{%
\subsection{Axis 3. Relational Mode (R, navigator Ř,
𝓕₄)}\label{axis-3.-relational-mode-r-navigator-ux159-ux1d4d5ux2084}}

Conventional meaning: direction and duality of the binding relation.

\begin{itemize}
\tightlist
\item
  \textbf{𐑾 {[}LR\_DUAL{]}, apex.} Fragment
  \texttt{lr⇔(x,\ y)\ ∧\ Θ(x,\ y)\ ∧\ ¬\ Θ(y,\ x)}.

  \begin{itemize}
  \tightlist
  \item
    FO: \texttt{(x\ ⇔\ y)\ ∧\ Θ(x,\ y)\ ∧\ ¬Θ(y,\ x)}; lateral
    biconditional binding with an asymmetric twist \texttt{Θ}.
  \item
    Named: a \textbf{left-right dual pair} with directed obstruction; a
    symmetric peer relation carrying a nonsymmetric cocycle (the bowtie
    duality of two sides that mutually determine yet do not dominate
    each other).
  \end{itemize}
\item
  \textbf{𐑽, close.} Fragment \texttt{f\ ⊣\ g\ ∧\ L\ Adj(f,\ g)}.

  \begin{itemize}
  \tightlist
  \item
    FO: \texttt{f\ ⊣\ g}; \texttt{f} is left adjoint to \texttt{g}.
  \item
    Named: an \textbf{adjunction} (unit/counit with the triangle
    identities); the categorical dagger-adjacent relation
    \texttt{A\ ⊣\ A†}.
  \end{itemize}
\item
  \textbf{𐑑, distant.} Fragment
  \texttt{Fun(x,\ y)\ ∧\ Nat(y,\ z)\ →\ Fun(x,\ z)}.

  \begin{itemize}
  \tightlist
  \item
    FO: \texttt{functor(x,\ y)\ ∧\ nat(y,\ z)\ →\ functor(x,\ z)}.
  \item
    Named: \textbf{functorial composition closed under natural
    transformation}; the 2-categorical composition law.
  \end{itemize}
\item
  \textbf{𐑩, distant.} Fragment \texttt{x\ ↑\ y\ ∧\ ¬(y\ ↑\ x)}.

  \begin{itemize}
  \tightlist
  \item
    FO: \texttt{dominates(x,\ y)\ ∧\ ¬dominates(y,\ x)}.
  \item
    Named: a \textbf{strict order / one-way supervisory relation}
    (hierarchy).
  \end{itemize}
\end{itemize}

\begin{center}\rule{0.5\linewidth}{0.5pt}\end{center}

\hypertarget{axis-4.-parity-symmetry-p-navigator-ux3c6-ux1d4d5ux2085}{%
\subsection{Axis 4. Parity / Symmetry (P, navigator Φ,
𝓕₅)}\label{axis-4.-parity-symmetry-p-navigator-ux3c6-ux1d4d5ux2085}}

Conventional meaning: the symmetry group fixing the object.

\begin{itemize}
\tightlist
\item
  \textbf{𐑹 {[}PM\_Z2{]}, apex.} Fragment
  \texttt{ℤ₂(x)\ ∧\ ∀g∈G(gx\ =\ x)\ ∧\ μ∘δ\ =\ id}.

  \begin{itemize}
  \tightlist
  \item
    FO:
    \texttt{ℤ₂\ ↷\ x\ ∧\ (∀g\ ∈\ G)(g·x\ =\ x)\ ∧\ (∀a)(μ(δ(a))\ =\ a)}.
  \item
    Named: a \textbf{split ℤ₂-retraction}: a ℤ₂ action under which
    \texttt{x} is invariant, together with a section/retraction pair
    \texttt{(δ,\ μ)} with \texttt{μ∘δ\ =\ id}. The signed subset-sum map
    is the canonical instance: \texttt{δ} distributes an element over
    \{in, out\}, \texttt{μ} sums, and their composite is the identity on
    the index set. This is the special Frobenius / Cartesian-bicategory
    identity, \texttt{μ} a split epi and \texttt{δ} a split mono.
  \end{itemize}
\item
  \textbf{𐑿, close.} Fragment
  \texttt{\textbar{}ψ⟩\ =\ Σ\ c\_i\ \textbar{}e\_i⟩}.

  \begin{itemize}
  \tightlist
  \item
    FO: \texttt{state(ψ)\ ∧\ ψ\ =\ Σ\_i\ c\_i\ e\_i} over an orthonormal
    basis.
  \item
    Named: a \textbf{superposition in a complex Hilbert space} (U(1)
    phase symmetry).
  \end{itemize}
\item
  \textbf{𐑬, close.} Fragment \texttt{ℤ₂(x)\ ∧\ ¬(x\ =\ -x)}.

  \begin{itemize}
  \tightlist
  \item
    FO: \texttt{ℤ₂\ ↷\ x\ ∧\ x\ ≠\ -x}.
  \item
    Named: a \textbf{nontrivial ℤ₂ representation} (sign flip with no
    fixed nonzero vector); the sign character acting freely.
  \end{itemize}
\item
  \textbf{𐑯, distant.} Fragment \texttt{∀g∈G(gx\ =\ x)}.

  \begin{itemize}
  \tightlist
  \item
    FO: \texttt{(∀g\ ∈\ G)(g·x\ =\ x)}.
  \item
    Named: a \textbf{G-invariant / fixed point of the group action}
    (full stabilizer).
  \end{itemize}
\item
  \textbf{𐑗, distant.} Fragment \texttt{¬∃sym(x)}.

  \begin{itemize}
  \tightlist
  \item
    FO: \texttt{¬∃σ(σ\ ≠\ id\ ∧\ σ·x\ =\ x)}.
  \item
    Named: an \textbf{asymmetric object} (trivial automorphism group).
  \end{itemize}
\end{itemize}

\begin{center}\rule{0.5\linewidth}{0.5pt}\end{center}

\hypertarget{axis-5.-fidelity-f-navigator-ux192-ux1d4d5ux2083}{%
\subsection{Axis 5. Fidelity (F, navigator ƒ,
𝓕₃)}\label{axis-5.-fidelity-f-navigator-ux192-ux1d4d5ux2083}}

Conventional meaning: information loss of the read/write map.

\begin{itemize}
\tightlist
\item
  \textbf{𐑐, apex.} Fragment \texttt{ℏ(x)\ ∧\ {[}x,\ p{]}\ =\ iℏ}.

  \begin{itemize}
  \tightlist
  \item
    FO: \texttt{quantum(x)\ ∧\ {[}x,\ p{]}\ =\ iℏ}.
  \item
    Named: a \textbf{canonical commutation relation}; lossless quantum
    channel (unitary, information preserving).
  \end{itemize}
\item
  \textbf{𐑞, close.} Fragment
  \texttt{Tr(ρ²)\ \textless{}\ 1\ ∧\ ρ\ =\ Σ\ p\_i\ \textbar{}i⟩⟨i\textbar{}}.

  \begin{itemize}
  \tightlist
  \item
    FO:
    \texttt{Tr(ρ²)\ \textless{}\ 1\ ∧\ ρ\ =\ Σ\_i\ p\_i\ \textbar{}i⟩⟨i\textbar{}}.
  \item
    Named: a \textbf{mixed state / density operator with nonzero
    entropy} (threshold fidelity, partial loss).
  \end{itemize}
\item
  \textbf{𐑱, distant.} Fragment \texttt{P(x)\ ∈\ \{0,1\}\ ∧\ det(x)}.

  \begin{itemize}
  \tightlist
  \item
    FO: \texttt{P(x)\ ∈\ \{0,\ 1\}\ ∧\ deterministic(x)}.
  \item
    Named: a \textbf{classical deterministic map} (projection to a
    definite outcome, lossy read).
  \end{itemize}
\end{itemize}

\begin{center}\rule{0.5\linewidth}{0.5pt}\end{center}

\hypertarget{axis-6.-kinetic-character-k-navigator-uxe7-ux1d4d5ux2085}{%
\subsection{Axis 6. Kinetic Character (K, navigator Ç,
𝓕₅)}\label{axis-6.-kinetic-character-k-navigator-uxe7-ux1d4d5ux2085}}

Conventional meaning: \textbf{potential}, how much the object still
holds, and whether the gate is open for it to go. τ is how long a
disturbance is retained against the observation window, so it reads as a
store, not a speed: τ ≪ T is potential already spent, τ ≫ T is potential
still held.

The ladder settles the reading on its own. If Ç were a rate, more τ
would be more of whatever Ç measures and the apex would sit at the
maximum. It does not. The apex is 𐑧 at ordinal 3, while 𐑪 (ordinal 4)
and 𐑺 (4.5) are both further out, and both are the ones carrying τ = ∞.
The peak is not the longest τ; it is the longest τ \textbf{whose gate is
open}, which is why \texttt{gate\_open(x)} appears in the apex fragment
and in none of the others. Past the apex the potential is still held and
can no longer move. A store that cannot discharge is not a larger store.

Potential alone does not move anything: Ç says how much is held, Ħ (Axis
10) says whether there is a direction to hold it toward. See
\emph{Reading the table} §4.

\begin{itemize}
\tightlist
\item
  \textbf{𐑧, apex.} Fragment
  \texttt{τ\ ≫\ T\ ∧\ eq(x)\ ∧\ gate\_open(x)}.

  \begin{itemize}
  \tightlist
  \item
    FO: \texttt{τ\ ≫\ T\ ∧\ equilibrium(x)\ ∧\ open(x)}.
  \item
    Named: a \textbf{quasi-static / adiabatic regime}: retained against
    the probe and still able to discharge. Maximal \emph{available}
    potential: the store is full and the gate is open. Conventionally, a
    reversible path: work is extractable precisely because relaxation is
    slow and nothing is pinned.
  \end{itemize}
\item
  \textbf{𐑤, close.} Fragment \texttt{τ\ ∼\ T\ ∧\ noisy(x)}.

  \begin{itemize}
  \tightlist
  \item
    FO: \texttt{τ\ ≈\ T\ ∧\ noisy(x)}.
  \item
    Named: a \textbf{critically damped / noisy regime} (store and probe
    on the same timescale; potential is leaking as fast as it is read).
  \end{itemize}
\item
  \textbf{𐑘, distant.} Fragment \texttt{τ\ ≪\ T\ ∧\ ∂\_t\ x\ =\ f(x)}.

  \begin{itemize}
  \tightlist
  \item
    FO: \texttt{τ\ ≪\ T\ ∧\ ẋ\ =\ f(x)}.
  \item
    Named: a \textbf{fast deterministic flow} (autonomous ODE). Nothing
    is retained: the potential discharges inside the observation window,
    so the object is always already at the bottom of its well.
  \end{itemize}
\item
  \textbf{𐑪, distant.} Fragment \texttt{τ\ =\ ∞\ ∧\ ord(x)}.

  \begin{itemize}
  \tightlist
  \item
    FO: \texttt{τ\ =\ ∞\ ∧\ ordered(x)}.
  \item
    Named: \textbf{frozen by order} (kinetically trapped, non-ergodic).
    The store is infinite and inaccessible: held forever, ordered, with
    no open gate. Maximum τ, unusable.
  \end{itemize}
\item
  \textbf{𐑺, distant.} Fragment \texttt{τ\ =\ ∞\ ∧\ dis(x)\ ∧\ MBL}.

  \begin{itemize}
  \tightlist
  \item
    FO: \texttt{τ\ =\ ∞\ ∧\ disordered(x)\ ∧\ MBL(x)}.
  \item
    Named: \textbf{many-body localization} (frozen by disorder,
    ergodicity broken). The same trap reached through disorder rather
    than order, which is why it ranks past 𐑪 rather than beside it.
  \end{itemize}
\end{itemize}

\begin{center}\rule{0.5\linewidth}{0.5pt}\end{center}

\hypertarget{axis-7.-scope-granularity-g-navigator-ux3b3-ux1d4d5ux2083}{%
\subsection{Axis 7. Scope / Granularity (G, navigator Γ,
𝓕₃)}\label{axis-7.-scope-granularity-g-navigator-ux3b3-ux1d4d5ux2083}}

Conventional meaning: subset-size regime of correlations.

\begin{itemize}
\tightlist
\item
  \textbf{𐑲, apex.} Fragment
  \texttt{∀y(y\ ⊂\ x\ →\ \textbar{}y\textbar{}\ \textless{}\ \textbar{}x\textbar{})}.

  \begin{itemize}
  \tightlist
  \item
    FO:
    \texttt{(∀y)(y\ ⊊\ x\ →\ \textbar{}y\textbar{}\ \textless{}\ \textbar{}x\textbar{})}.
  \item
    Named: a \textbf{Dedekind-finite object}: no proper subset is
    equinumerous with the whole. Strict cardinal drop under proper
    inclusion (fine-grained, global correlations resolved at every
    scale).
  \end{itemize}
\item
  \textbf{𐑔, close.} Fragment
  \texttt{∃y∈x(\textbar{}y\textbar{}\ ∼\ \textbar{}x\textbar{})}.

  \begin{itemize}
  \tightlist
  \item
    FO:
    \texttt{(∃y\ ⊊\ x)(\textbar{}y\textbar{}\ =\ \textbar{}x\textbar{})}.
  \item
    Named: a \textbf{Dedekind-infinite object}: a proper subset in
    bijection with the whole (intermediate / collective scale,
    self-similarity present).
  \end{itemize}
\item
  \textbf{𐑚, distant.} Fragment
  \texttt{∀y∈x(\textbar{}y\textbar{}\ \textless{}\ \textbar{}x\textbar{})}.

  \begin{itemize}
  \tightlist
  \item
    FO:
    \texttt{(∀y\ ∈\ x)(\textbar{}y\textbar{}\ \textless{}\ \textbar{}x\textbar{})}.
  \item
    Named: \textbf{element-wise strict smallness} (each member below the
    whole; local / mesoscale, no member captures the aggregate).
  \end{itemize}
\end{itemize}

\begin{center}\rule{0.5\linewidth}{0.5pt}\end{center}

\hypertarget{axis-8.-interaction-grammar-ux3b3-navigator-ux262-ux1d4d5ux2084}{%
\subsection{Axis 8. Interaction Grammar (Γ, navigator ɢ,
𝓕₄)}\label{axis-8.-interaction-grammar-ux3b3-navigator-ux262-ux1d4d5ux2084}}

Conventional meaning: the logical connective of the coupling.

\begin{itemize}
\tightlist
\item
  \textbf{𐑵 {[}BROADCAST\_TRANSCENDENCE{]}, apex.} Fragment
  \texttt{f\ →\ all(x)\ ∧\ broadcast(x,\ f)}.

  \begin{itemize}
  \tightlist
  \item
    FO: \texttt{(∀z)(f(z))\ ∧\ broadcast(x,\ f)}; one map applied to
    all, one-to-all coupling.
  \item
    Named: a \textbf{universal / broadcast quantifier} (\texttt{∀} as a
    single coupling to the whole domain); a transcendence atom,
    exceeding finitary composition.
  \end{itemize}
\item
  \textbf{𐑠 {[}SEQAX{]}, close.} Fragment
  \texttt{seq!(f,\ g)\ ∧\ ⟨→⟩(f,\ g,\ τ)\ ∧\ ¬\ ⟨→⟩(g,\ f,\ τ)}.

  \begin{itemize}
  \tightlist
  \item
    FO: \texttt{seq(f,\ g)\ ∧\ (f\ →\_τ\ g)\ ∧\ ¬(g\ →\_τ\ f)}.
  \item
    Named: a \textbf{strict sequencing / causal order} (\texttt{f} then
    \texttt{g}, not the reverse); the sequential-composition axiom,
    temporal directedness.
  \end{itemize}
\item
  \textbf{𐑜, distant.} Fragment \texttt{f\ ∨\ g\ ∨\ h}.

  \begin{itemize}
  \tightlist
  \item
    FO: \texttt{f\ ∨\ g\ ∨\ h}.
  \item
    Named: \textbf{disjunctive coupling} (any condition suffices;
    parallel alternatives).
  \end{itemize}
\item
  \textbf{𐑝, distant.} Fragment \texttt{f\ ∧\ g\ ∧\ h}.

  \begin{itemize}
  \tightlist
  \item
    FO: \texttt{f\ ∧\ g\ ∧\ h}.
  \item
    Named: \textbf{conjunctive coupling} (all conditions required;
    simultaneous).
  \end{itemize}
\end{itemize}

\begin{center}\rule{0.5\linewidth}{0.5pt}\end{center}

\hypertarget{axis-9.-criticality-ux3c6-navigator-ux1d4d5ux2085}{%
\subsection{Axis 9. Criticality (Φ, navigator ⊙,
𝓕₅)}\label{axis-9.-criticality-ux3c6-navigator-ux1d4d5ux2085}}

Conventional meaning: analytic character of the fixed point.

\begin{itemize}
\tightlist
\item
  \textbf{⊙ {[}PHI\_C{]}, apex.} Fragment
  \texttt{ξ\ →\ ∞\ ∧\ μ∘δ\ =\ id}.

  \begin{itemize}
  \tightlist
  \item
    FO: \texttt{ξ\ →\ ∞\ ∧\ (∀a)(μ(δ(a))\ =\ a)}; correlation length
    diverges while the split/fuse pair remains an exact identity.
  \item
    Named: the \textbf{critical fixed point with exact Frobenius
    closure}: a scale-invariant point (\texttt{ξ\ →\ ∞}) at which the
    section/retraction is lossless. The Frobenius section survives the
    divergence, the self-dual critical point.
  \end{itemize}
\item
  \textbf{𐑮, close.} Fragment \texttt{ξ\ ∈\ ℂ\ ∧\ Im(ξ)\ →\ ∞}.

  \begin{itemize}
  \tightlist
  \item
    FO: \texttt{ξ\ ∈\ ℂ\ ∧\ Im(ξ)\ →\ ∞}.
  \item
    Named: a \textbf{complex-axis criticality} (Lee-Yang edge, complex
    RG fixed point, ζ-zero line); analytic continuation off the real
    axis.
  \end{itemize}
\item
  \textbf{𐑻, distant.} Fragment
  \texttt{H(λ)\ non-Herm\ ∧\ det(H\ -\ λI)\ =\ 0\ ∧\ ∂\_λ\ H\ =\ 0}.

  \begin{itemize}
  \tightlist
  \item
    FO:
    \texttt{nonHermitian(H(λ))\ ∧\ det(H\ -\ λI)\ =\ 0\ ∧\ ∂\_λ\ H\ =\ 0}.
  \item
    Named: a \textbf{non-Hermitian exceptional point}: eigenvalue and
    eigenvector coalescence, a square-root branch point of the spectrum.
  \end{itemize}
\item
  \textbf{𐑣, distant.} Fragment \texttt{ξ\ →\ ∞\ ∧\ chaotic(x)}.

  \begin{itemize}
  \tightlist
  \item
    FO: \texttt{ξ\ →\ ∞\ ∧\ chaotic(x)}.
  \item
    Named: \textbf{divergence into chaos} (unbounded correlation without
    a fixed point; supercritical runaway).
  \end{itemize}
\item
  \textbf{𐑢, distant.} Fragment \texttt{¬∃ξ(diverges(ξ))}.

  \begin{itemize}
  \tightlist
  \item
    FO: \texttt{¬∃ξ(diverges(ξ))}.
  \item
    Named: a \textbf{subcritical / gapped phase} (no divergent length
    scale; stable).
  \end{itemize}
\end{itemize}

\begin{center}\rule{0.5\linewidth}{0.5pt}\end{center}

\hypertarget{axis-10.-chirality-h-navigator-ux127-ux1d4d5ux2084}{%
\subsection{Axis 10. Chirality (H, navigator Ħ,
𝓕₄)}\label{axis-10.-chirality-h-navigator-ux127-ux1d4d5ux2084}}

Conventional meaning: handedness of the state under mirror and shift,
which is to say the direction available to whatever Ç holds. H is
chirality throughout, never temporal depth or memory.

The ladder is a tower and its descent. 𐑫 is a tower of
\texttt{μ∘δ}-fixed points at every rank, standing, with nothing falling
through it. 𐑖 is a rank descent: something is falling. 𐑓 is
shift-invariance, no handedness at all, so nothing \emph{can} fall. Read
with Axis 6: Ç is the store, Ħ is the direction it discharges into, and
neither moves anything alone. See \emph{Reading the table} §4.

\begin{itemize}
\tightlist
\item
  \textbf{𐑫 {[}ETERNAL\_FIXEDPOINT{]}, apex.} Fragment
  \texttt{∀n∃φ(rank(φ)\ \textgreater{}\ n\ ∧\ φ\ fixed\ by\ μ∘δ\ ∧\ φ\ ∈\ V)}.

  \begin{itemize}
  \tightlist
  \item
    FO:
    \texttt{(∀n)(∃φ)(rank(φ)\ \textgreater{}\ n\ ∧\ μ(δ(φ))\ =\ φ\ ∧\ φ\ ∈\ V)}.
  \item
    Named: an \textbf{inexhaustible tower of Frobenius fixed points} of
    unbounded rank; a topologically protected chirality with fixed
    points at every height (a proper class of \texttt{μ∘δ}-fixed
    objects).
  \end{itemize}
\item
  \textbf{𐑖 {[}TEMPD2{]}, close.} Fragment
  \texttt{∃y∃z(y\ ∈\ x\ ∧\ z\ ∈\ y\ ∧\ ¬\ z\ ∈\ x\ ∧\ rank(z)\ \textless{}\ rank(y))}.

  \begin{itemize}
  \tightlist
  \item
    FO: as written; a membership descent that is not transitive
    (\texttt{z\ ∈\ y\ ∈\ x} but \texttt{z\ ∉\ x}).
  \item
    Named: a \textbf{non-transitive membership chain / depth-2
    well-founded descent}; handedness that survives two levels of
    descent. Note what the fragment does:
    \texttt{rank(z)\ \textless{}\ rank(y)} is a fall through rank. The
    tower of 𐑫 is standing; here it is spilling.
  \end{itemize}
\item
  \textbf{𐑒, distant.} Fragment \texttt{∃y(P(y)\ ↔\ P(S²(y)))}.

  \begin{itemize}
  \tightlist
  \item
    FO: \texttt{(∃y)(P(y)\ ↔\ P(S²(y)))}.
  \item
    Named: a \textbf{period-2 symmetry under the shift} (invariance
    under \texttt{S²}, soft chirality).
  \end{itemize}
\item
  \textbf{𐑓, distant.} Fragment \texttt{∀x(P(x)\ ↔\ P(S(x)))}.

  \begin{itemize}
  \tightlist
  \item
    FO: \texttt{(∀x)(P(x)\ ↔\ P(S(x)))}.
  \item
    Named: \textbf{full shift invariance} (period-1; achiral, no
    handedness to break).
  \end{itemize}
\end{itemize}

\begin{center}\rule{0.5\linewidth}{0.5pt}\end{center}

\hypertarget{axis-11.-stoichiometry-s-navigator-ux3c3-ux1d4d5ux2083}{%
\subsection{Axis 11. Stoichiometry (S, navigator Σ,
𝓕₃)}\label{axis-11.-stoichiometry-s-navigator-ux3c3-ux1d4d5ux2083}}

Conventional meaning: matching cardinality of the two coupled sides A
and B.

\begin{itemize}
\tightlist
\item
  \textbf{𐑳, apex.} Fragment \texttt{∃a∈A∃b∈B(type(a)\ ≠\ type(b))}.

  \begin{itemize}
  \tightlist
  \item
    FO: \texttt{(∃a\ ∈\ A)(∃b\ ∈\ B)(type(a)\ ≠\ type(b))}.
  \item
    Named: a \textbf{heterotypic / unmatched coupling} (n:m, the two
    sides carry distinct types; cross-type reaction center present).
  \end{itemize}
\item
  \textbf{𐑕, close.} Fragment \texttt{∀a∈A∀b∈B(type(a)\ =\ type(b))}.

  \begin{itemize}
  \tightlist
  \item
    FO: \texttt{(∀a\ ∈\ A)(∀b\ ∈\ B)(type(a)\ =\ type(b))}.
  \item
    Named: a \textbf{homotypic matched coupling} (n:n, all pairs share a
    type).
  \end{itemize}
\item
  \textbf{𐑙, distant.} Fragment
  \texttt{\textbar{}A\textbar{}\ =\ 1\ ∧\ \textbar{}B\textbar{}\ =\ 1}.

  \begin{itemize}
  \tightlist
  \item
    FO:
    \texttt{\textbar{}A\textbar{}\ =\ 1\ ∧\ \textbar{}B\textbar{}\ =\ 1}.
  \item
    Named: a \textbf{1:1 singleton pairing} (a single bond, the minimal
    stoichiometric unit).
  \end{itemize}
\end{itemize}

\begin{center}\rule{0.5\linewidth}{0.5pt}\end{center}

\hypertarget{axis-12.-topological-protection-ux3c9-navigator-ux3c9-ux1d4d5ux2084}{%
\subsection{Axis 12. Topological Protection (Ω, navigator Ω,
𝓕₄)}\label{axis-12.-topological-protection-ux3c9-navigator-ux3c9-ux1d4d5ux2084}}

Conventional meaning: the homotopy invariant guarding the state.

\begin{itemize}
\tightlist
\item
  \textbf{𐑟 {[}BRAID\_TRANSCENDENCE{]}, apex.} Fragment
  \texttt{Braid(σ\_i)\ ∧\ R\_matrix\ ≠\ 0\ ∧\ nonAbelian(x)}.

  \begin{itemize}
  \tightlist
  \item
    FO: \texttt{braid(σ\_i)\ ∧\ R\ ≠\ 0\ ∧\ nonAbelian(x)}.
  \item
    Named: a \textbf{non-Abelian braiding / representation of the braid
    group} with a nontrivial R-matrix (anyonic protection); a
    transcendence atom.
  \end{itemize}
\item
  \textbf{𐑭 {[}ZWIND{]}, close.} Fragment
  \texttt{∮\_γ\ A\ =\ 2πn\ ∧\ n\ ∈\ ℤ\ ∧\ wind(γ)\ ≠\ 0}.

  \begin{itemize}
  \tightlist
  \item
    FO: \texttt{∮\_γ\ A\ =\ 2πn\ ∧\ n\ ∈\ ℤ\ ∧\ winding(γ)\ ≠\ 0}.
  \item
    Named: an \textbf{integer winding number / first Chern class}; ℤ
    topological protection (quantized holonomy).
  \end{itemize}
\item
  \textbf{𐑴, distant.} Fragment \texttt{∮\_γ\ A\ =\ nπ\ ∧\ n\ ∈\ ℤ₂}.

  \begin{itemize}
  \tightlist
  \item
    FO: \texttt{∮\_γ\ A\ =\ nπ\ ∧\ n\ ∈\ ℤ₂}.
  \item
    Named: a \textbf{ℤ₂ Berry phase} (half-integer holonomy, ℤ₂
    protection).
  \end{itemize}
\item
  \textbf{𐑷, distant.} Fragment \texttt{∮\_γ\ dx\ =\ 0}.

  \begin{itemize}
  \tightlist
  \item
    FO: \texttt{∮\_γ\ dx\ =\ 0}.
  \item
    Named: an \textbf{exact / trivial holonomy} (no protection;
    contractible loop).
  \end{itemize}
\end{itemize}

\begin{center}\rule{0.5\linewidth}{0.5pt}\end{center}

\hypertarget{reading-the-table}{%
\subsection{Reading the table}\label{reading-the-table}}

Four structural facts carry into the paper.

\begin{enumerate}
\def\labelenumi{\arabic{enumi}.}
\item
  \textbf{The apex value of each axis is where a named transcendence or
  Frobenius condition appears.} \texttt{μ∘δ\ =\ id} recurs at the apex
  of Parity (P 𐑹), Criticality (Φ ⊙) and Chirality (H 𐑫), and L8's two
  declared transcendence atoms (BROADCAST at Γ 𐑵, BRAID at Ω 𐑟) sit at
  axis apices. The apex column is not decorative: it is exactly the set
  of values that exceed the ZFC\_fe baseline \textbf{on the road to L8}.
  That qualifier is load bearing, see 3.
\item
  \textbf{The twelve axes are independent conventional invariants.}
  Dimension, connectivity, adjunction direction, symmetry group, channel
  fidelity, relaxation regime, cardinal granularity, logical connective,
  criticality type, chirality, stoichiometric matching, and homotopy
  class are twelve quantities a mathematician already computes
  separately. The crystal asserts they form a complete coordinate system
  for structural type. Phase 2 argues the order in which they must be
  fixed is the magnum-opus stage sequence.
\item
  \textbf{Transcendence is not a single summit, and L9 proves it.} CLINK
  L8 called itself the terminal layer of the chain. It is not: CLINK L9,
  the Gaussian Moat Resolution layer, sits above it at O\_∞⁺, reached by
  eight promotions the navigator notes as HODGE BRIDGE TRANSCENDENCE.
  The striking part is the direction. L9 exceeds L8 by
  \textbf{relinquishing both of L8's transcendence atoms}: Γ gives up
  BROADCAST 𐑵 for 𐑝 (\texttt{f\ ∧\ g\ ∧\ h}, a three-unit stitch) at the
  maximum gap of 1.0, and Ω gives up BRAID 𐑟 for 𐑭 (integer winding). Ð
  drops from the imscriptive fixed point 𐑦 to 𐑛, which this table calls
  \texttt{distant} and reads as a 0-dimensional finite set, and which L9
  names PRIME\_POINT: a point-like prime atom, exactly what a Gaussian
  prime is. Four axes hold fixed across the rung, and they are the ones
  that matter: ƒ 𐑐, ⊙ ⊙, H 𐑫, S 𐑳. The self-modeling criticality and the
  eternal Frobenius fixed point are precisely what L9 does not give up.
\item
  \textbf{Ç holds, Ħ points, and the pair is what moves.} The two axes
  are read together or not at all. Ç is a store: how much the object
  still has, and whether its gate is open. Ħ is a direction: whether
  there is anywhere for the store to go. Neither moves anything alone. A
  full store with no handedness (Ç 𐑪, Ħ 𐑓) is frozen: τ = ∞ and
  shift-invariant, so nothing can fall. A direction with an empty store
  (Ç 𐑘, Ħ 𐑫) is a standing tower with nothing descending it.

  The fragments carry this without needing to be told. Ħ's apex 𐑫 is a
  tower of \texttt{μ∘δ}-fixed points at every rank; its next value 𐑖 is
  \texttt{rank(z)\ \textless{}\ rank(y)}, a fall through rank. So 𐑫 is
  the tower standing and 𐑖 is the tower spilling, and the step from 𐑫
  down to 𐑖 is not a loss but a discharge. This is also the cleanest
  evidence for the potential reading of Ç: on a rate reading, an axis's
  apex should be its longest τ, and Ç's apex is 𐑧 at ordinal 3 while
  both τ = ∞ values rank further out. Only \texttt{gate\_open(x)}
  explains that, and only a store has a gate.

  Both facts are formal, not glossed.
  \texttt{Smaragdine\_Synthetica/lean} and the p4rakernel tree carry the
  pair as theorems: the bare tick is Frobenius-minimal in eleven
  coordinates and maximal in chirality alone (an empty store, an endless
  tower), while lived duration sits at Ç's open gate with chirality one
  step down (a full store, descending). The gap between the two is the ⊙
  seam, which is exactly where a store is permitted to become a flow.

  So the apex column is a coordinate in L8's dialect, not a ladder to a
  summit. Read a value as low because it is \texttt{distant} here and
  you will call L9 a demotion; measured in its own register it is a tier
  above. The Frobenius conditions carry across dialects. The
  transcendence atoms do not.
\end{enumerate}

\end{document}
